\section{Un modèle gaussien à deux facteurs pour le taux OAT}
\label{sec:modele}

\subsection{Dynamique et moments gaussiens}

Le taux court sans risque s'écrit $r(t) = \alpha(t) + x(t)$ et le spread court OAT–€STR
$s(t) = \beta(t) + y(t)$, où $\alpha$ et $\beta$ sont déterministes et où $x$ et $y$ sont deux processus
d'Ornstein-Uhlenbeck centrés \citep[éq.~4.4-4.5]{BrigoMercurio2006} :
\begin{equation}
\dd x = -a_x\,x\,\dd t + \sigma_x\,\dd W_x,\qquad \dd y = -a_y\,y\,\dd t + \sigma_y\,\dd W_y,\qquad
\dd W_x\,\dd W_y = \rho\,\dd t,
\label{eq:dynamique}
\end{equation}
avec $x(0) = y(0) = 0$.
Le taux court de l'OAT est $\bar r = r + s = \bar\alpha + x + y$, avec $\bar\alpha = \alpha + \beta$. Dans les
notations de \citet{BrigoMercurio2006}, $(a_x, a_y, \sigma_x, \sigma_y)$ sont $(a, b, \sigma, \eta)$ et
$\bar\alpha$ est $\varphi$.

Toutes les formules du modèle se ramènent aux moments, conditionnels à $\F_t$, du vecteur gaussien formé de
$x(T)$, $y(T)$, $\int_t^T x$ et $\int_t^T y$. Avec $\tau = T - t$, la duration stochastique d'un zéro-coupon est
\begin{equation}
H_a(t,T) = \frac{1-e^{-a\tau}}{a}, \qquad \E\Big[\int_t^T x(u)\,\dd u \,\Big|\, \F_t\Big] = H_{a_x}(t,T)\,x(t),
\end{equation}
et la variance de l'intégrale d'un facteur de paramètres $(a,\sigma)$ vaut
\begin{equation}
V_{a,\sigma}(t,T) = \Var\Big[\int_t^T x(u)\,\dd u\Big]
= \frac{\sigma^2}{a^2}\Big[\tau + \frac{2}{a}e^{-a\tau} - \frac{1}{2a}e^{-2a\tau} - \frac{3}{2a}\Big].
\end{equation}
Le terme croisé est deux fois la covariance des intégrales \citep[lemme~4.2.1, éq.~4.10]{BrigoMercurio2006} :
\begin{equation}
V_{xy}(t,T) = \frac{2\rho\sigma_x\sigma_y}{a_x a_y}\Big[\tau - H_{a_x}(t,T) - H_{a_y}(t,T)
+ \frac{1-e^{-(a_x+a_y)\tau}}{a_x+a_y}\Big],
\end{equation}
et la variance totale $\bar V = V_x + V_y + V_{xy} = \Var\big[\int_t^T (x+y)\big]$ est la variance $V(t,T)$ du
livre. Les variances et covariances des facteurs eux-mêmes, qui servent au pricing, sont données en annexe~A.

\subsection{Le facteur taux seul : Hull-White et Jamshidian}

Lorsque $\sigma_y = 0$, le modèle se réduit à celui de Hull et White sous la forme $r = \alpha + x$
\citep[section~3.3]{BrigoMercurio2006}. Le terme déterministe recale la courbe €STR :
$\int_0^T \alpha = -\log P^M(0,T) + \tfrac12 V_x(0,T)$. Comme $\int_t^T x$ est gaussien de moyenne
$H_{a_x}(t,T)\,x(t)$ et de variance $V_x(t,T)$, le prix d'un zéro-coupon est log-affine en $x(t)$
\citep[éq.~3.39]{BrigoMercurio2006} :
\begin{equation}
P(t,T) = \frac{P^M(0,T)}{P^M(0,t)}\,\exp\Big(-H_{a_x}(t,T)\,x(t) + \tfrac12\big[V_x(t,T) - V_x(0,T) + V_x(0,t)\big]\Big).
\label{eq:zc_hw}
\end{equation}

Une option européenne d'échéance $T_0$ et de strike $X$ sur une obligation de flux $K_i$ aux dates $T_i$ a, dans
ce modèle, un prix exact. Le prix de l'obligation en $T_0$ étant décroissant en $x(T_0)$, il existe un unique
$x^\ast$ tel que $\sum_i K_i P(T_0,T_i;x^\ast) = X$, et l'option se décompose en une somme d'options sur
zéro-coupons de strikes $X_i = P(T_0,T_i;x^\ast)$, chacune de formule fermée \citep{Jamshidian1989} ;
\citet[section~3.11.1]{BrigoMercurio2006} et \citet[section~6.3.6]{Chaix2021} en donnent la mise en œuvre. Une
swaption payeuse est un put de strike 1 sur l'obligation formée de la jambe fixe et du nominal : elle se price de
la même façon \citep[éq.~3.44-3.46]{BrigoMercurio2006}. À la monnaie, son prix de marché en volatilité normale
$\sigma_N$ est $A\,\sigma_N\sqrt{T_0/2\pi}$, où $A$ est l'annuité du swap sous-jacent.

\subsection{Le spread comme second facteur}

\citet{Russo} lisent les deux facteurs du G2++ comme le taux sans risque et le spread de crédit. On modélise ici
directement le spread court OAT–€STR. Il contient la prime de crédit de l'État français, mais aussi des primes de
liquidité et de collatéral : il est traité comme un facteur de risque de marché, sans chercher à en isoler la part
de défaut.

Les modèles de crédit à intensité justifient la forme $\bar r = r + y$. Supposons que le défaut soit un état
absorbant, qu'il ne cause pas les facteurs de taux, lesquels ne sautent pas au défaut, et que l'événement de
défaut lui-même ne soit pas pricé. Avec un recouvrement en valeur de marché, où le porteur récupère au défaut une
fraction $\zeta$ du prix qu'aurait eu le titre sans défaut, le prix avant défaut d'un zéro-coupon risqué est
\begin{equation}
\bar P(t,T) = \E^{\Q}_t\Big[e^{-\int_t^T (r_u + \tilde\lambda_u)\,\dd u}\Big],\qquad \tilde\lambda \approx (1-\zeta)\,\lambda,
\end{equation}
où $\lambda$ est l'intensité de défaut et $\tilde\lambda$ l'intensité ajustée du recouvrement
\citep[cours~5, hypothèses~II.1-II.3 et prop.~II.1-II.2]{DuffieSingleton1999,MPR2025}. Les prix ne
révèlent que le produit de la perte et de l'intensité, jamais l'une sans l'autre, et une prime de liquidité
s'ajoute de la même façon au taux d'actualisation (éq.~II.15 du cours). \citet[section~22.4]{BrigoMercurio2006}
écrivent la même forme pour un recouvrement nul et une intensité corrélée aux taux. Le spread total $s = \beta + y$
est donc modélisé par un facteur gaussien sans paramètre de perte : $y$ est une intensité ajustée qui agrège
crédit, liquidité et collatéral. Ce cadre fonde la forme du taux d'actualisation avant défaut ; le mémoire ne
modélise ni la date de défaut ni le saut du prix au défaut.

Le terme $\bar\alpha$ recale directement la courbe OAT \citep[corollaire~4.2.1, éq.~4.13]{BrigoMercurio2006} :
\begin{equation}
\int_0^T\bar\alpha(u)\,\dd u = -\log\bar P^M(0,T) + \tfrac12\bar V(0,T),
\label{eq:recalage}
\end{equation}
et la différence $\bar\alpha - \alpha = \beta$ est le spread forward de marché augmenté des termes de convexité du
spread et du terme croisé. Comme $\int_t^T(x+y)$ est gaussien de moyenne $H_{a_x}x(t) + H_{a_y}y(t)$ et de
variance $\bar V(t,T)$, le zéro-coupon OAT vaut \citep[théorème~4.2.1]{BrigoMercurio2006}
\begin{equation}
\bar P(t,T) = \frac{\bar P^M(0,T)}{\bar P^M(0,t)}\exp\Big(-H_{a_x}(t,T)\,x(t) - H_{a_y}(t,T)\,y(t)
+ \tfrac12\big[\bar V(t,T) - \bar V(0,T) + \bar V(0,t)\big]\Big).
\label{eq:zc_oat}
\end{equation}
C'est l'expression des équations~16 à~24 de \citet{Russo}. Une OAT de flux $K_i$ aux dates $T_i$ vaut
$\bar B(t) = \sum_i K_i \bar P(t,T_i)$ (prix dirty), et ses durations stochastiques sont
$D_x = \sum_i w_i H_{a_x}(t,T_i)$ et $D_y = \sum_i w_i H_{a_y}(t,T_i)$, avec $w_i = K_i\bar P(t,T_i)/\bar B(t)$.
Lorsque les deux vitesses de retour sont faibles, $H_{a_x} \approx H_{a_y}$ : les deux facteurs déplacent alors
les prix des OAT de la même façon, propriété qui reviendra à la section~\ref{sec:resultats}.
